% Source: https://artofproblemsolving.com/community/c3h2520854_expected_value_problem_
% Stable identity: aops-topic-2520854-q01
If she does not try, exactly $1$ person dies. If she tries, $0$ people die with probability $p$ and $2$ die with probability $1-p$, so the expected number of deaths is $2(1-p)$.

\textbf{(a)} $p=\dfrac12$: $2\cdot\dfrac12=1$, the same as not trying. Both choices give $1$ expected death, so in terms of expected deaths it makes no difference.

\textbf{(b)} $p=0.4$: $2\cdot0.6=1.2>1$, so she should \textbf{not} try.

\textbf{(c)} Trying is worse exactly when $2(1-p)>1$, i.e. $p<\dfrac12$. So the odds are against trying for success rates in the interval $\boxed{\left(0,\dfrac12\right)}$; for $p>\dfrac12$ trying is better.
